\documentclass[12pt,a4paper]{article}
\usepackage[utf8]{vietnam}
\usepackage{amsmath,amssymb}
\usepackage{tabularx}
\usepackage{tikz}
\usepackage{xcolor}
\usepackage[top=0.6cm, bottom=0.6cm, left=0.6cm, right=0.6cm]{geometry}
\usepackage[dethi]{ex_test}

\SAOPTN{widthSA=2.6,heightSA=0.62}
\setlength{\parskip}{0pt}

% Dong ke cham mo chuan viet tay
\newcommand{\dongke}[1]{%
  \begingroup%
  \setlength{\parskip}{0pt}%
  \baselineskip=18.5pt%
  \vspace{1.0mm}%
  \foreach \i in {1,...,#1}{%
    \noindent{\color{gray!60}\dotfill}\par%
  }%
  \vspace{0.6mm}%
  \endgroup%
}

% Dinh nghia o bai tap tu luan trong bang
\newcommand{\oBaiTap}[3][6]{%
  \vspace{1.0mm}%
  \noindent\textbf{#2} #3\par%
  \dongke{#1}%
}

\begin{document}
\pagestyle{empty}
% =============================================================================
% ========================= MA DE 101 - MAT TRUOC (TRANG 1/2) =============
% =============================================================================
\noindent
\begin{tabularx}{\linewidth}{|X|c|}
\hline
\textbf{PHIẾU KIỂM TRA ĐỊNH KỲ TOÁN 8 - TUẦN 07 - BUỔI 02} & \textbf{Mã đề: 101 - Trang 1/2} \\
\hline
\end{tabularx}

\vspace{1.2mm}

\noindent
\begin{tabularx}{\linewidth}{|p{3.2cm}|X|}
\hline
\centering\textbf{Điểm số} & \centering\textbf{Lời nhận xét của thầy cô} \tabularnewline
\rule{0pt}{1.8cm} & \\
\hline
\end{tabularx}

\vspace{1.2mm}

\noindent\textbf{Câu 1 (0,25 điểm).} Kết quả của $2x(3x - 1)$ là: \shortans{$6x^2 - 2x$}\par\vspace{0.3mm}
\noindent\textbf{Câu 2 (0,25 điểm).} Bậc của đa thức $M = 3x^2 y^2 - 2x^3 + x^2 y^2 - 1$ là: \shortans{$4$}\par\vspace{0.3mm}
\noindent\textbf{Câu 3 (0,25 điểm).} Giá trị của $A = 2x^2 y + 3xy$ tại $x = 1$, $y = 2$ là: \shortans{$10$}\par\vspace{0.3mm}
\noindent\textbf{Câu 4 (0,25 điểm).} Khai triển $(x + 5)^2$ ta được $x^2 + \ldots + 25$. Điền vào chỗ trống: \shortans{$10x$}\par\vspace{0.3mm}
\noindent\textbf{Câu 5 (0,25 điểm).} Phân tích $4x^2 - 9$ thành nhân tử được: \shortans{$(2x-3)(2x+3)$}\par\vspace{0.3mm}
\noindent\textbf{Câu 6 (0,25 điểm).} Phân tích đa thức $xy + x + y + 1$ thành nhân tử được: \shortans{$(x+1)(y+1)$}\par\vspace{0.3mm}
\noindent\textbf{Câu 7 (0,25 điểm).} Hình bình hành là tứ giác có các cạnh đối như thế nào? \shortans{song song}\par\vspace{0.3mm}
\noindent\textbf{Câu 8 (0,25 điểm).} Cho hình bình hành $ABCD$ có $\widehat{A} = 110^\circ$. Số đo góc $\widehat{C}$ là: \shortans{$110^\circ$}\par\vspace{0.3mm}
\noindent\textbf{Câu 9 (0,25 điểm).} Cho hình bình hành $ABCD$ có hai đường chéo cắt nhau tại $O$, biết $AC = 12\text{ cm}$. Độ dài đoạn $OA$ là: \shortans{$6\text{ cm}$}\par\vspace{0.3mm}
\noindent\textbf{Câu 10 (0,25 điểm).} Một mảnh vườn hình bình hành có độ dài hai cạnh kề là $15\text{ m}$ và $20\text{ m}$. Chu vi mảnh vườn là: \shortans{$70\text{ m}$}\par\vspace{0.3mm}

\vspace{0.5mm}
\noindent\textbf{Câu 11 (2,0 điểm). Thực hiện phép tính (hợp lý nếu có thể):}
\vspace{0.5mm}

\noindent
\begin{tabularx}{\linewidth}{|X|X|}
\hline
\oBaiTap[10]{a)}{$3x^2(2x - 1) - 2x(3x^2 - 4x)$} & \oBaiTap[10]{b)}{$(x + 2)^2 - (x - 1)^2$} \\
\hline
\oBaiTap[10]{c)}{$(20x^4 y^3 - 15x^3 y^2 + 5x^2 y) : 5x^2 y$} & \oBaiTap[10]{d)}{$(x - 1)(x + 4) - x(x + 2)$} \\
\hline
\end{tabularx}

\newpage
% =============================================================================
% ========================= MA DE 101 - MAT SAU (TRANG 2/2) ==============
% =============================================================================
\noindent
\begin{tabularx}{\linewidth}{|X|}
\hline
\centering\textbf{TOÁN 8 - TUẦN 07 - BUỔI 02 | Mã đề: 101 - Trang 2/2} \\
\hline
\end{tabularx}

\vspace{1.2mm}

\noindent\textbf{Câu 12 (2,0 điểm). Tìm $x$, biết:}
\vspace{0.5mm}

\noindent
\begin{tabularx}{\linewidth}{|X|X|}
\hline
\oBaiTap[5]{a)}{$2x(x - 5) - 2x^2 = 20$} & \oBaiTap[5]{b)}{$x(x - 7) + 3(x - 7) = 0$} \\
\hline
\oBaiTap[5]{c)}{$x^2 - 9x = 0$} & \oBaiTap[5]{d)}{$(x - 3)^2 - 16 = 0$} \\
\hline
\end{tabularx}

\vspace{1.2mm}

\noindent
\begin{tabularx}{\linewidth}{|X|}
\hline
\vspace{0.5mm}
\textbf{Câu 13 (2,5 điểm).} Cho hình bình hành $ABCD$. Gọi $M, N$ lần lượt là trung điểm của các cạnh $AB$ và $CD$.\par
\textbf{a) (1,25 điểm)} Chứng minh tứ giác $AMCN$ là hình bình hành.\par
\textbf{b) (1,25 điểm)} Gọi $O$ là giao điểm của hai đường chéo $AC$ và $BD$. Chứng minh ba điểm $M, O, N$ thẳng hàng.\par
\vspace{0.4mm}
\noindent\textit{Bài giải:}\dongke{11}
\\ \hline
\end{tabularx}

\vspace{1.2mm}

\noindent
\begin{tabularx}{\linewidth}{|X|}
\hline
\vspace{0.5mm}
\textbf{Câu 14 (1,0 điểm).} Cho biểu thức $A = x^2 - 8x + 2035$. Tìm giá trị nhỏ nhất của biểu thức $A$.\par
\vspace{0.4mm}
\noindent\textit{Bài giải:}\dongke{8}
\\ \hline
\end{tabularx}

\newpage
% =============================================================================
% ========================= MA DE 102 - MAT TRUOC (TRANG 1/2) =============
% =============================================================================
\noindent
\begin{tabularx}{\linewidth}{|X|c|}
\hline
\textbf{PHIẾU KIỂM TRA ĐỊNH KỲ TOÁN 8 - TUẦN 07 - BUỔI 02} & \textbf{Mã đề: 102 - Trang 1/2} \\
\hline
\end{tabularx}

\vspace{1.2mm}

\noindent
\begin{tabularx}{\linewidth}{|p{3.2cm}|X|}
\hline
\centering\textbf{Điểm số} & \centering\textbf{Lời nhận xét của thầy cô} \tabularnewline
\rule{0pt}{1.8cm} & \\
\hline
\end{tabularx}

\vspace{1.2mm}

\noindent\textbf{Câu 1 (0,25 điểm).} Kết quả của $3x(2x - 3)$ là: \shortans{$6x^2 - 9x$}\par\vspace{0.3mm}
\noindent\textbf{Câu 2 (0,25 điểm).} Bậc của đa thức $N = 4x^3 y - 3x^2 + 2x^3 y + 5$ là: \shortans{$4$}\par\vspace{0.3mm}
\noindent\textbf{Câu 3 (0,25 điểm).} Giá trị của $B = 3x^2 y + 2xy$ tại $x = 1$, $y = 3$ là: \shortans{$15$}\par\vspace{0.3mm}
\noindent\textbf{Câu 4 (0,25 điểm).} Khai triển $(x + 6)^2$ ta được $x^2 + \ldots + 36$. Điền vào chỗ trống: \shortans{$12x$}\par\vspace{0.3mm}
\noindent\textbf{Câu 5 (0,25 điểm).} Phân tích $9x^2 - 16$ thành nhân tử được: \shortans{$(3x-4)(3x+4)$}\par\vspace{0.3mm}
\noindent\textbf{Câu 6 (0,25 điểm).} Phân tích đa thức $xy + 2x + y + 2$ thành nhân tử được: \shortans{$(x+1)(y+2)$}\par\vspace{0.3mm}
\noindent\textbf{Câu 7 (0,25 điểm).} Tứ giác có các cạnh đối song song gọi là hình gì? \shortans{hình bình hành}\par\vspace{0.3mm}
\noindent\textbf{Câu 8 (0,25 điểm).} Cho hình bình hành $ABCD$ có $\widehat{A} = 120^\circ$. Số đo góc $\widehat{C}$ là: \shortans{$120^\circ$}\par\vspace{0.3mm}
\noindent\textbf{Câu 9 (0,25 điểm).} Cho hình bình hành $ABCD$ có hai đường chéo cắt nhau tại $O$, biết $BD = 16\text{ cm}$. Độ dài đoạn $OB$ là: \shortans{$8\text{ cm}$}\par\vspace{0.3mm}
\noindent\textbf{Câu 10 (0,25 điểm).} Một mảnh vườn hình bình hành có độ dài hai cạnh kề là $12\text{ m}$ và $18\text{ m}$. Chu vi mảnh vườn là: \shortans{$60\text{ m}$}\par\vspace{0.3mm}

\vspace{0.5mm}
\noindent\textbf{Câu 11 (2,0 điểm). Thực hiện phép tính (hợp lý nếu có thể):}
\vspace{0.5mm}

\noindent
\begin{tabularx}{\linewidth}{|X|X|}
\hline
\oBaiTap[10]{a)}{$2x^2(3x - 2) - 3x(2x^2 - 3x)$} & \oBaiTap[10]{b)}{$(x + 3)^2 - (x - 2)^2$} \\
\hline
\oBaiTap[10]{c)}{$(24x^4 y^3 - 18x^3 y^2 + 6x^2 y) : 6x^2 y$} & \oBaiTap[10]{d)}{$(x - 2)(x + 5) - x(x + 1)$} \\
\hline
\end{tabularx}

\newpage
% =============================================================================
% ========================= MA DE 102 - MAT SAU (TRANG 2/2) ==============
% =============================================================================
\noindent
\begin{tabularx}{\linewidth}{|X|}
\hline
\centering\textbf{TOÁN 8 - TUẦN 07 - BUỔI 02 | Mã đề: 102 - Trang 2/2} \\
\hline
\end{tabularx}

\vspace{1.2mm}

\noindent\textbf{Câu 12 (2,0 điểm). Tìm $x$, biết:}
\vspace{0.5mm}

\noindent
\begin{tabularx}{\linewidth}{|X|X|}
\hline
\oBaiTap[5]{a)}{$3x(x - 4) - 3x^2 = 24$} & \oBaiTap[5]{b)}{$x(x - 6) + 4(x - 6) = 0$} \\
\hline
\oBaiTap[5]{c)}{$x^2 - 8x = 0$} & \oBaiTap[5]{d)}{$(x - 4)^2 - 25 = 0$} \\
\hline
\end{tabularx}

\vspace{1.2mm}

\noindent
\begin{tabularx}{\linewidth}{|X|}
\hline
\vspace{0.5mm}
\textbf{Câu 13 (2,5 điểm).} Cho hình bình hành $ABCD$. Gọi $E, F$ lần lượt là trung điểm của các cạnh $AB$ và $CD$.\par
\textbf{a) (1,25 điểm)} Chứng minh tứ giác $AECF$ là hình bình hành.\par
\textbf{b) (1,25 điểm)} Gọi $O$ là giao điểm của hai đường chéo $AC$ và $BD$. Chứng minh ba điểm $E, O, F$ thẳng hàng.\par
\vspace{0.4mm}
\noindent\textit{Bài giải:}\dongke{11}
\\ \hline
\end{tabularx}

\vspace{1.2mm}

\noindent
\begin{tabularx}{\linewidth}{|X|}
\hline
\vspace{0.5mm}
\textbf{Câu 14 (1,0 điểm).} Cho biểu thức $B = x^2 - 6x + 2031$. Tìm giá trị nhỏ nhất của biểu thức $B$.\par
\vspace{0.4mm}
\noindent\textit{Bài giải:}\dongke{8}
\\ \hline
\end{tabularx}

\newpage
% =============================================================================
% ========================= MA DE 103 - MAT TRUOC (TRANG 1/2) =============
% =============================================================================
\noindent
\begin{tabularx}{\linewidth}{|X|c|}
\hline
\textbf{PHIẾU KIỂM TRA ĐỊNH KỲ TOÁN 8 - TUẦN 07 - BUỔI 02} & \textbf{Mã đề: 103 - Trang 1/2} \\
\hline
\end{tabularx}

\vspace{1.2mm}

\noindent
\begin{tabularx}{\linewidth}{|p{3.2cm}|X|}
\hline
\centering\textbf{Điểm số} & \centering\textbf{Lời nhận xét của thầy cô} \tabularnewline
\rule{0pt}{1.8cm} & \\
\hline
\end{tabularx}

\vspace{1.2mm}

\noindent\textbf{Câu 1 (0,25 điểm).} Kết quả của $4x(2x - 1)$ là: \shortans{$8x^2 - 4x$}\par\vspace{0.3mm}
\noindent\textbf{Câu 2 (0,25 điểm).} Bậc của đa thức $P = 5x^2 y^3 - 4x^4 + x^2 y^3 - 2$ là: \shortans{$5$}\par\vspace{0.3mm}
\noindent\textbf{Câu 3 (0,25 điểm).} Giá trị của $C = 4x^2 y - 2xy$ tại $x = 1$, $y = 2$ là: \shortans{$4$}\par\vspace{0.3mm}
\noindent\textbf{Câu 4 (0,25 điểm).} Khai triển $(x + 7)^2$ ta được $x^2 + \ldots + 49$. Điền vào chỗ trống: \shortans{$14x$}\par\vspace{0.3mm}
\noindent\textbf{Câu 5 (0,25 điểm).} Phân tích $16x^2 - 25$ thành nhân tử được: \shortans{$(4x-5)(4x+5)$}\par\vspace{0.3mm}
\noindent\textbf{Câu 6 (0,25 điểm).} Phân tích đa thức $xy + 3x + y + 3$ thành nhân tử được: \shortans{$(x+1)(y+3)$}\par\vspace{0.3mm}
\noindent\textbf{Câu 7 (0,25 điểm).} Tứ giác có các cạnh đối song song và bằng nhau là hình gì? \shortans{hình bình hành}\par\vspace{0.3mm}
\noindent\textbf{Câu 8 (0,25 điểm).} Cho hình bình hành $ABCD$ có $\widehat{A} = 105^\circ$. Số đo góc $\widehat{C}$ là: \shortans{$105^\circ$}\par\vspace{0.3mm}
\noindent\textbf{Câu 9 (0,25 điểm).} Cho hình bình hành $ABCD$ có hai đường chéo cắt nhau tại $O$, biết $AC = 14\text{ cm}$. Độ dài đoạn $OC$ là: \shortans{$7\text{ cm}$}\par\vspace{0.3mm}
\noindent\textbf{Câu 10 (0,25 điểm).} Một mảnh vườn hình bình hành có độ dài hai cạnh kề là $16\text{ m}$ và $24\text{ m}$. Chu vi mảnh vườn là: \shortans{$80\text{ m}$}\par\vspace{0.3mm}

\vspace{0.5mm}
\noindent\textbf{Câu 11 (2,0 điểm). Thực hiện phép tính (hợp lý nếu có thể):}
\vspace{0.5mm}

\noindent
\begin{tabularx}{\linewidth}{|X|X|}
\hline
\oBaiTap[10]{a)}{$4x^2(2x - 1) - 2x(4x^2 - 5x)$} & \oBaiTap[10]{b)}{$(x + 4)^2 - (x - 1)^2$} \\
\hline
\oBaiTap[10]{c)}{$(28x^4 y^3 - 21x^3 y^2 + 7x^2 y) : 7x^2 y$} & \oBaiTap[10]{d)}{$(x - 3)(x + 6) - x(x + 2)$} \\
\hline
\end{tabularx}

\newpage
% =============================================================================
% ========================= MA DE 103 - MAT SAU (TRANG 2/2) ==============
% =============================================================================
\noindent
\begin{tabularx}{\linewidth}{|X|}
\hline
\centering\textbf{TOÁN 8 - TUẦN 07 - BUỔI 02 | Mã đề: 103 - Trang 2/2} \\
\hline
\end{tabularx}

\vspace{1.2mm}

\noindent\textbf{Câu 12 (2,0 điểm). Tìm $x$, biết:}
\vspace{0.5mm}

\noindent
\begin{tabularx}{\linewidth}{|X|X|}
\hline
\oBaiTap[5]{a)}{$4x(x - 3) - 4x^2 = 24$} & \oBaiTap[5]{b)}{$x(x - 8) + 2(x - 8) = 0$} \\
\hline
\oBaiTap[5]{c)}{$x^2 - 7x = 0$} & \oBaiTap[5]{d)}{$(x - 2)^2 - 36 = 0$} \\
\hline
\end{tabularx}

\vspace{1.2mm}

\noindent
\begin{tabularx}{\linewidth}{|X|}
\hline
\vspace{0.5mm}
\textbf{Câu 13 (2,5 điểm).} Cho hình bình hành $MNPQ$. Gọi $E, F$ lần lượt là trung điểm của các cạnh $MN$ và $PQ$.\par
\textbf{a) (1,25 điểm)} Chứng minh tứ giác $MEPF$ là hình bình hành.\par
\textbf{b) (1,25 điểm)} Gọi $I$ là giao điểm của hai đường chéo $MP$ và $NQ$. Chứng minh ba điểm $E, I, F$ thẳng hàng.\par
\vspace{0.4mm}
\noindent\textit{Bài giải:}\dongke{11}
\\ \hline
\end{tabularx}

\vspace{1.2mm}

\noindent
\begin{tabularx}{\linewidth}{|X|}
\hline
\vspace{0.5mm}
\textbf{Câu 14 (1,0 điểm).} Cho biểu thức $C = x^2 - 10x + 2048$. Tìm giá trị nhỏ nhất của biểu thức $C$.\par
\vspace{0.4mm}
\noindent\textit{Bài giải:}\dongke{8}
\\ \hline
\end{tabularx}

\newpage
% =============================================================================
% ========================= MA DE 104 - MAT TRUOC (TRANG 1/2) =============
% =============================================================================
\noindent
\begin{tabularx}{\linewidth}{|X|c|}
\hline
\textbf{PHIẾU KIỂM TRA ĐỊNH KỲ TOÁN 8 - TUẦN 07 - BUỔI 02} & \textbf{Mã đề: 104 - Trang 1/2} \\
\hline
\end{tabularx}

\vspace{1.2mm}

\noindent
\begin{tabularx}{\linewidth}{|p{3.2cm}|X|}
\hline
\centering\textbf{Điểm số} & \centering\textbf{Lời nhận xét của thầy cô} \tabularnewline
\rule{0pt}{1.8cm} & \\
\hline
\end{tabularx}

\vspace{1.2mm}

\noindent\textbf{Câu 1 (0,25 điểm).} Kết quả của $5x(2x - 3)$ là: \shortans{$10x^2 - 15x$}\par\vspace{0.3mm}
\noindent\textbf{Câu 2 (0,25 điểm).} Bậc của đa thức $Q = 2x^4 y - 3x^3 + 3x^4 y + 7$ là: \shortans{$5$}\par\vspace{0.3mm}
\noindent\textbf{Câu 3 (0,25 điểm).} Giá trị của $D = 2x^2 y - 5xy$ tại $x = 1$, $y = -2$ là: \shortans{$6$}\par\vspace{0.3mm}
\noindent\textbf{Câu 4 (0,25 điểm).} Khai triển $(x + 8)^2$ ta được $x^2 + \ldots + 64$. Điền vào chỗ trống: \shortans{$16x$}\par\vspace{0.3mm}
\noindent\textbf{Câu 5 (0,25 điểm).} Phân tích $25x^2 - 36$ thành nhân tử được: \shortans{$(5x-6)(5x+6)$}\par\vspace{0.3mm}
\noindent\textbf{Câu 6 (0,25 điểm).} Phân tích đa thức $xy + 4x + y + 4$ thành nhân tử được: \shortans{$(x+1)(y+4)$}\par\vspace{0.3mm}
\noindent\textbf{Câu 7 (0,25 điểm).} Tứ giác có hai đường chéo cắt nhau tại trung điểm mỗi đường là hình gì? \shortans{hình bình hành}\par\vspace{0.3mm}
\noindent\textbf{Câu 8 (0,25 điểm).} Cho hình bình hành $ABCD$ có $\widehat{A} = 115^\circ$. Số đo góc $\widehat{C}$ là: \shortans{$115^\circ$}\par\vspace{0.3mm}
\noindent\textbf{Câu 9 (0,25 điểm).} Cho hình bình hành $ABCD$ có hai đường chéo cắt nhau tại $O$, biết $BD = 18\text{ cm}$. Độ dài đoạn $OD$ là: \shortans{$9\text{ cm}$}\par\vspace{0.3mm}
\noindent\textbf{Câu 10 (0,25 điểm).} Một mảnh vườn hình bình hành có độ dài hai cạnh kề là $25\text{ m}$ và $35\text{ m}$. Chu vi mảnh vườn là: \shortans{$120\text{ m}$}\par\vspace{0.3mm}

\vspace{0.5mm}
\noindent\textbf{Câu 11 (2,0 điểm). Thực hiện phép tính (hợp lý nếu có thể):}
\vspace{0.5mm}

\noindent
\begin{tabularx}{\linewidth}{|X|X|}
\hline
\oBaiTap[10]{a)}{$3x^2(3x - 2) - x(9x^2 - 8x)$} & \oBaiTap[10]{b)}{$(x + 5)^2 - (x - 2)^2$} \\
\hline
\oBaiTap[10]{c)}{$(32x^4 y^3 - 24x^3 y^2 + 8x^2 y) : 8x^2 y$} & \oBaiTap[10]{d)}{$(x - 1)(x + 6) - x(x + 3)$} \\
\hline
\end{tabularx}

\newpage
% =============================================================================
% ========================= MA DE 104 - MAT SAU (TRANG 2/2) ==============
% =============================================================================
\noindent
\begin{tabularx}{\linewidth}{|X|}
\hline
\centering\textbf{TOÁN 8 - TUẦN 07 - BUỔI 02 | Mã đề: 104 - Trang 2/2} \\
\hline
\end{tabularx}

\vspace{1.2mm}

\noindent\textbf{Câu 12 (2,0 điểm). Tìm $x$, biết:}
\vspace{0.5mm}

\noindent
\begin{tabularx}{\linewidth}{|X|X|}
\hline
\oBaiTap[5]{a)}{$5x(x - 2) - 5x^2 = 30$} & \oBaiTap[5]{b)}{$x(x - 5) + 3(x - 5) = 0$} \\
\hline
\oBaiTap[5]{c)}{$x^2 - 6x = 0$} & \oBaiTap[5]{d)}{$(x - 5)^2 - 49 = 0$} \\
\hline
\end{tabularx}

\vspace{1.2mm}

\noindent
\begin{tabularx}{\linewidth}{|X|}
\hline
\vspace{0.5mm}
\textbf{Câu 13 (2,5 điểm).} Cho hình bình hành $ABCD$. Gọi $I, K$ lần lượt là trung điểm của các cạnh $AB$ và $CD$.\par
\textbf{a) (1,25 điểm)} Chứng minh tứ giác $AICK$ là hình bình hành.\par
\textbf{b) (1,25 điểm)} Gọi $O$ là giao điểm của hai đường chéo $AC$ và $BD$. Chứng minh ba điểm $I, O, K$ thẳng hàng.\par
\vspace{0.4mm}
\noindent\textit{Bài giải:}\dongke{11}
\\ \hline
\end{tabularx}

\vspace{1.2mm}

\noindent
\begin{tabularx}{\linewidth}{|X|}
\hline
\vspace{0.5mm}
\textbf{Câu 14 (1,0 điểm).} Cho biểu thức $D = x^2 - 4x + 2028$. Tìm giá trị nhỏ nhất của biểu thức $D$.\par
\vspace{0.4mm}
\noindent\textit{Bài giải:}\dongke{8}
\\ \hline
\end{tabularx}

\newpage
\end{document}
